\chapter{Transformer Architecture Deep Dive}
\label{chap:transformers}

The Transformer architecture, introduced by Vaswani et al. in "Attention Is All You Need" \cite{vaswani2017attention}, marked a paradigm shift in natural language processing. Prior to 2017, Recurrent Neural Networks (RNNs) and Long Short-Term Memory networks (LSTMs) \cite{hochreiter1997long} dominated sequence modeling. However, RNNs processed data sequentially, making parallelization impossible and struggling with long-range dependencies.

Transformers solved both problems by discarding recurrence entirely in favor of an attention mechanism that processes the entire sequence at once.

\section{The Core Mechanism: Self-Attention}

At its core, self-attention allows a model to weigh the importance of different words in a sentence relative to each other. When processing the word "bank" in "river bank", the model should pay attention to "river" to disambiguate meaning.

\subsection{Query, Key, and Value}
Self-attention is calculated using three vectors derived from each input token's embedding:
\begin{itemize}
    \item \textbf{Query ($Q$):} What the current token is looking for.
    \item \textbf{Key ($K$):} What the current token offers to others.
    \item \textbf{Value ($V$):} The actual content of the token.
\end{itemize}

These vectors are obtained by multiplying the input embedding $X$ by learned weight matrices $W^Q$, $W^K$, and $W^V$.

\subsection{Scaled Dot-Product Attention}
The attention score is computed by taking the dot product of the Query with all Keys. A high dot product indicates high relevance.
\begin{equation}
    \text{Attention}(Q, K, V) = \text{softmax}\left(\frac{QK^T}{\sqrt{d_k}}\right)V
\end{equation}
\begin{itemize}
    \item $QK^T$: The similarity matrix (how much focus each word places on every other word).
    \item $\sqrt{d_k}$: A scaling factor to prevent gradients from vanishing in the softmax function.
    \item Softmax: Normalizes the scores so they sum to 1.
    \item $V$: The values are weighted by the attention scores.
\end{itemize}

\section{Multi-Head Attention}

A single attention head might focus on syntactic relationships (e.g., subject-verb agreement). Another might focus on semantic relationships. To capture multiple types of relationships simultaneously, Transformers use Multi-Head Attention.

\begin{equation}
    \text{MultiHead}(Q, K, V) = \text{Concat}(\text{head}_1, \dots, \text{head}_h)W^O
\end{equation}
where each $\text{head}_i = \text{Attention}(QW_i^Q, KW_i^K, VW_i^V)$.

This allows the model to attend to information from different representation subspaces at different positions.

\section{Positional Encodings}

Since Transformers process all tokens in parallel, they have no inherent sense of order. "The cat ate the mouse" looks the same as "The mouse ate the cat" to a raw attention mechanism.

To fix this, we inject positional information into the input embeddings. The original paper used sine and cosine functions of different frequencies:
\begin{align}
    PE_{(pos, 2i)} &= \sin(pos / 10000^{2i/d_{model}}) \\
    PE_{(pos, 2i+1)} &= \cos(pos / 10000^{2i/d_{model}})
\end{align}
Modern models like Llama 2 \cite{touvron2023llama} use Rotary Positional Embeddings (RoPE), which rotate the query and key vectors in complex space to encode relative positions.

\section{The Architecture Blocks}

A standard Transformer block consists of two sub-layers:
\begin{enumerate}
    \item Multi-Head Self-Attention
    \item Position-wise Feed-Forward Network (FFN)
\end{enumerate}
Each sub-layer has a residual connection followed by layer normalization:
\begin{equation}
    \text{LayerNorm}(x + \text{Sublayer}(x))
\end{equation}

\subsection{Feed-Forward Network}
The FFN is applied to each position separately and identically. It consists of two linear transformations with a ReLU activation in between:
\begin{equation}
    \text{FFN}(x) = \max(0, xW_1 + b_1)W_2 + b_2
\end{equation}
This is where the model stores its "knowledge" or facts, while the attention mechanism handles "reasoning" or routing information.

\section{Encoder vs. Decoder Architectures}

The original Transformer had both an Encoder (to process input) and a Decoder (to generate output). Over time, architectures diverged.

\subsection{Encoder-Only (BERT)}
Models like BERT \cite{devlin2018bert} use only the Encoder stack. They are \textbf{bidirectional}, meaning they can see the entire context (left and right) at once.
\begin{itemize}
    \item \textbf{Use Case:} Classification, Named Entity Recognition, Sentiment Analysis.
    \item \textbf{Training Objective:} Masked Language Modeling (predicting missing words).
\end{itemize}

\subsection{Decoder-Only (GPT)}
Models like GPT-3 \cite{brown2020language} and Llama \cite{touvron2023llama} use only the Decoder stack. They are \textbf{autoregressive}, meaning they generate one token at a time, attending only to previous tokens.
\begin{itemize}
    \item \textbf{Use Case:} Text Generation, Chatbots, Code Completion.
    \item \textbf{Training Objective:} Causal Language Modeling (predicting the next word).
\end{itemize}

\subsection{Encoder-Decoder (T5, BART)}
These models keep both stacks. The encoder processes the input text, and the decoder generates the output text \cite{raffel2020exploring}.
\begin{itemize}
    \item \textbf{Use Case:} Translation, Summarization.
\end{itemize}

\section{Vision Transformers (ViT)}

The Transformer architecture is not limited to text. The Vision Transformer (ViT) \cite{dosovitskiy2020image} splits an image into fixed-size patches (e.g., 16x16 pixels), linearly embeds each patch, adds position embeddings, and feeds the sequence of vectors to a standard Transformer encoder. This has largely replaced CNNs in state-of-the-art computer vision.

\section{Conclusion}

The Transformer's ability to scale with compute and data has led to the current AI boom. Its simple building blocks---attention and feed-forward networks---combined with massive datasets, allow it to learn intricate patterns in human language, code, and even biology (e.g., AlphaFold).

In the next chapter, we will discuss how to align these powerful models with human intent using Reinforcement Learning.
